% answers_short.tex —— 1990 年·广东卷·文史类 · 简短答案速查
% 由 tools/make_short_answers.py 自 answers.tex 的 \gksolhead/\gkanswer 提取，
% 并与 1990/1990_广东_答案要点.md 逐题比对后生成，勿手改（改源 MD 后重跑）。

\begin{gknotice}
\textbf{使用说明}：本节只给最终答案便于快速核对，完整推导见解析版。答案为固定来源 PDF 卷内答案表与解答所载（出版者/版次未核实），非官方标准答案；第二卷第三至五题卷内只注「见理工农医类解答」，另附项目推导。
\end{gknotice}

\gkpart{选择题}
\begin{center}\begin{tabular}{cccccccccccccc}
\toprule
题号 & 1 & 2 & 3 & 4 & 5 & 6 & 7 & 8 & 9 & 10 & 11 & 12 & 13 \\
\midrule
答案 & A & D & A & D & B & B & A & A & B & C & C & C & D \\
\bottomrule
\end{tabular}\end{center}\medskip
\begin{center}\begin{tabular}{ccccccccccccc}
\toprule
题号 & 14 & 15 & 16 & 17 & 18 & 19 & 20 & 21 & 22 & 23 & 24 & 25 \\
\midrule
答案 & B & D & C & D & B & C & B & A & C & B & C & D \\
\bottomrule
\end{tabular}\end{center}

\gkpart{填空题}
\begin{itemize}[leftmargin=2.6em]
  \item[\textbf{26.}] (1) $\dfrac{3\sqrt{10}}{10}$；(2) $-10$；(3) $\dfrac43$；(4) $\dfrac{\sqrt6}{3}a$；(5) $5$
\end{itemize}

\gkpart{解答题}
\begin{itemize}[leftmargin=2.6em]
  \item[\textbf{27.}] （1）定义域为 $\left\{x\ \middle|\ \dfrac12<x<4\right\}$；（2）不等式的解集为 $\left\{x\ \middle|\ \dfrac12<x<2\right\}$。
  \item[\textbf{28.}] 卷内未列结果，仅注「见理工农医类第二卷第二题解答」。（1）为证明题；（2）二面角为 $60^\circ$（此值据题面推得，非卷内所载，见解析）。
  \item[\textbf{29.}] 卷内未列结果，仅注「见理工农医类第二卷第三题解答」。（1）椭圆方程为 $\dfrac{x^2}{4}+\dfrac{y^2}{3}=1$；（2）$\cos\angle F_1PF_2=\dfrac35$（此二值据题面推得，非卷内所载，见解析）。
  \item[\textbf{30.}] 卷内未列结果，仅注「见理工农医类第二卷第四题解答」。（1）$a_1=\dfrac{\pi}{4}$，公差 $d=\dfrac{\pi}{3}$；（2）为证明题（结论均据题面推得，非卷内所载，见解析）。
\end{itemize}
